% SPDX-License-Identifier: LPPL-1.3c
% Copyright (C) 2026 Christophe Jorssen
% Author and Current Maintainer: Christophe Jorssen <christophe.jorssen@gmail.com>
% LPPL maintenance status: maintained
% This file is part of luafemm. See LICENSE and LICENSES.md for its terms.
\documentclass[varwidth=155mm,border=3mm]{standalone}
\usepackage{tikz,pgfplots}
\usetikzlibrary{femm}
\usepgfplotslibrary{femm}
\pgfplotsset{compat=1.18}
\newbox\BiasCalculation
\tikzset{femm/materials/core/.style={
  bh={0/0,.5/50,1/150,1.5/600,1.8/3000,2/10000}},
  femm/arrow spacing=18mm}
\newcommand\PaintImportedCore{%
  \fill[even odd rule,blue!7] (-13,-13) rectangle (13,13)
    (-7,-7) rectangle (7,7);
  \fill[orange!25] (-19,-6) rectangle (-15,6) (-5,-6) rectangle (-2,6);}
\begin{document}
\centering
% A producer: this could instead be a supported DC answer written by FEMM.
\setbox\BiasCalculation=\hbox{%
\begin{tikzpicture}[femm/problem={depth=20,interpolation=femm,
  xmin=-25,xmax=25,ymin=-20,ymax=20,mesh size=1.8},
  femm/circuit={name=drive,current=4},
  femm/export={format=ans,file=\jobname-bias.ans}]
  \path[even odd rule,femm/region={material=core}]
    (-13,-13) rectangle (13,13) (-7,-7) rectangle (7,7);
  \path[femm/region={material=copper,circuit=drive,turns=100}]
    (-19,-6) rectangle (-15,6);
  \path[femm/region={material=copper,circuit=drive,turns=-100}]
    (-5,-6) rectangle (-2,6);
  \femmsolve
\end{tikzpicture}}
% Restore, display, then capture an owned operating point at picture end.
\begin{tikzpicture}[scale=1.35,
  femm/import={format=ans,file=\jobname-bias.ans},
  femm/operating point={name=imported-core}]
  \PaintImportedCore
  \pic {femm geometry};
  \pic[femm/lines=13] {femm field};
  \path[femm/profile={name=bias-cut,samples=81}] (-12,10)--(12,10);
  \node[above,font=\small] at (0,20) {Imported DC bias, $I_0=4$ A};
\end{tikzpicture}\hspace{5mm}%
\begin{tikzpicture}[scale=1.35,
  femm/problem={operating point=imported-core,permeability=incremental,
    frequency=200,circuit currents={drive=.1}}]
  \PaintImportedCore
  \pic {femm geometry};
  \pic[femm/lines=13,femm/field style={red!65!black}] {femm field};
  \path[femm/profile={name=response-cut,samples=81}] (-12,10)--(12,10);
  \node[above,font=\small] at (0,20) {Response, $\widehat{\delta I}=0.1$ A};
\end{tikzpicture}
\par\medskip
\foreach \Profile/\View/\Title in
  {bias-cut/real/{DC $B_x$ (T)},response-cut/real/{Real $\delta B_x$ (T)}}{%
\begin{tikzpicture}
\begin{axis}[width=7.2cm,height=5.2cm,grid=major,
  xlabel={$x$ along the upper limb (mm)},ylabel={\Title}]
  \femmplot[profile=\Profile,abscissa=x,component=Bx,representation=\View,
    plot options={blue!70!black,thick,mark=none}]
\end{axis}
\end{tikzpicture}}
\par\smallskip\small
Illustrative single-valued steel law. The incremental field uses the unchanged
bias mesh and the differential reluctivity. Its contours have their own levels;
it is a small response at 200 Hz, not the total field or a new nonlinear bias.
\end{document}
